\documentclass[11pt,a4paper]{article}

%====================================================
% Encodage, langue, typographie
%====================================================
\usepackage[utf8]{inputenc}
\usepackage[T1]{fontenc}
\usepackage[french]{babel}
\usepackage{lmodern}
\usepackage{microtype}
\usepackage{setspace}

%====================================================
% Mise en page
%====================================================
\usepackage[a4paper,margin=2cm]{geometry}
\usepackage{fancyhdr}
\usepackage{lastpage}
\usepackage{titlesec}
\usepackage{enumitem}
\usepackage{array,tabularx,booktabs,multicol}
\usepackage{hyperref}
\hypersetup{colorlinks=true,linkcolor=blue!50!black,urlcolor=blue!60!black}

%====================================================
% Maths, chimie, unités
%====================================================
\usepackage{amsmath,amssymb,amsthm}
\usepackage{siunitx}
\sisetup{locale=FR,detect-all,per-mode=symbol}
\usepackage[version=4]{mhchem}
\usepackage{chemfig}

%====================================================
% Graphiques et encadrés
%====================================================
\usepackage{tikz}
\usetikzlibrary{arrows.meta,calc,positioning,decorations.pathmorphing,patterns}
\usepackage[most]{tcolorbox}

%====================================================
% Couleurs
%====================================================
\definecolor{bleuprepa}{RGB}{20,55,110}
\definecolor{vertprepa}{RGB}{20,105,70}
\definecolor{rougeprepa}{RGB}{160,35,35}
\definecolor{orangeprepa}{RGB}{190,100,20}
\definecolor{violetprepa}{RGB}{90,60,145}
\definecolor{grisclair}{RGB}{245,247,250}
\definecolor{bleuclair}{RGB}{235,243,255}
\definecolor{vertclair}{RGB}{235,250,240}
\definecolor{rougeclair}{RGB}{255,238,238}
\definecolor{orangeclair}{RGB}{255,246,232}
\definecolor{violetclair}{RGB}{246,241,255}

%====================================================
% En-têtes
%====================================================
\pagestyle{fancy}
\fancyhf{}
\fancyhead[L]{Première spé PC}
\fancyhead[C]{Constitution et transformations de la matière}
\fancyhead[R]{Page \thepage/\pageref{LastPage}}
\fancyfoot[C]{\small Cours complet - méthodes, exemples, exercices et corrigés}
\setlength{\headheight}{14pt}

%====================================================
% Titres
%====================================================
\titleformat{\section}{\Large\bfseries\color{bleuprepa}}{\thesection.}{0.6em}{}
\titleformat{\subsection}{\large\bfseries\color{violetprepa}}{\thesubsection.}{0.6em}{}
\titleformat{\subsubsection}{\bfseries\color{orangeprepa}}{\thesubsubsection.}{0.6em}{}

%====================================================
% Boîtes
%====================================================
\tcbset{enhanced,breakable,sharp corners=downhill,boxrule=0.7pt,colback=white,fonttitle=\bfseries,before skip=8pt,after skip=8pt}
\newtcolorbox{definitionbox}[1][]{colframe=bleuprepa,colback=bleuclair,title=Définition,#1}
\newtcolorbox{propositionbox}[1][]{colframe=vertprepa,colback=vertclair,title=Propriété,#1}
\newtcolorbox{theorembox}[1][]{colframe=violetprepa,colback=violetclair,title=Résultat fondamental,#1}
\newtcolorbox{methodbox}[1][]{colframe=bleuprepa,colback=white,title=Méthode,#1}
\newtcolorbox{retenirbox}[1][]{colframe=vertprepa,colback=vertclair,title=À retenir,#1}
\newtcolorbox{erreurbox}[1][]{colframe=rougeprepa,colback=rougeclair,title=Erreur classique,#1}
\newtcolorbox{exemplebox}[1][]{colframe=orangeprepa,colback=orangeclair,title=Exemple corrigé,#1}
\newtcolorbox{demobox}[1][]{colframe=black!60,colback=grisclair,title=Démonstration / justification,#1}
\newtcolorbox{exercicebox}[1][]{colframe=bleuprepa,colback=white,title=Exercice,#1}
\newtcolorbox{correctionbox}[1][]{colframe=black!65,colback=grisclair,title=Correction détaillée,#1}

%====================================================
% Commandes utiles
%====================================================
\newcommand{\M}{\mathrm{M}}
\newcommand{\n}{\mathrm{n}}
\newcommand{\m}{\mathrm{m}}
\newcommand{\V}{\mathrm{V}}
\newcommand{\cconc}{\mathrm{c}}
\newcommand{\xmax}{x_{\max}}
\newcommand{\xf}{x_f}
\newcommand{\molL}{\si{mol.L^{-1}}}
\newcommand{\gL}{\si{g.L^{-1}}}
\newcommand{\Na}{N_{\mathrm A}}
\newcommand{\E}{\mathrm{E}}
\newcommand{\Rdt}{\eta}

\setlength{\parindent}{0pt}
\setlength{\parskip}{4pt}
\onehalfspacing

\begin{document}

%====================================================
% Page de titre
%====================================================
\begin{titlepage}
\begin{center}
\vspace*{0.8cm}
{\Huge\bfseries Constitution et transformations de la matière}\\[0.25cm]
{\LARGE Cours complet de première générale}\\[0.25cm]
{\Large Spécialité physique-chimie}\\[0.8cm]

\begin{tikzpicture}[scale=1]
\node[draw=bleuprepa,fill=bleuclair,rounded corners,minimum width=13cm,minimum height=2cm] (a) {\Large \textbf{Mole, transformations, redox, titrages}};
\node[draw=vertprepa,fill=vertclair,rounded corners,minimum width=13cm,minimum height=2cm,below=0.45cm of a] (b) {\Large \textbf{Structure, polarité, cohésion, solubilité}};
\node[draw=orangeprepa,fill=orangeclair,rounded corners,minimum width=13cm,minimum height=2cm,below=0.45cm of b] (c) {\Large \textbf{Chimie organique, synthèses, combustions}};
\draw[-{Latex},thick,bleuprepa] (a) -- (b);
\draw[-{Latex},thick,bleuprepa] (b) -- (c);
\end{tikzpicture}

\vfill

\begin{tcolorbox}[colframe=bleuprepa,colback=bleuclair,width=0.92\textwidth,title=Objectif du document]
Ce polycopié présente un cours progressif, rigoureux et exploitable en classe de première spécialité physique-chimie. Il couvre les trois grands blocs du thème : suivi de transformation chimique, structure microscopique et propriétés de la matière, chimie organique, synthèses et combustions. Chaque partie contient des définitions, propriétés, méthodes, exemples corrigés, exercices progressifs et corrections détaillées.
\end{tcolorbox}

\vfill

\begin{tabular}{rl}
\textbf{Niveau :} & Première générale - spécialité physique-chimie.\\
\textbf{Public :} & élèves solides, préparation rigoureuse, support professeur.\\
\textbf{Contenu :} & cours, méthodes, erreurs classiques, exercices, devoir bilan, synthèse.\\
\textbf{Compilation :} & PDFLaTeX ou Overleaf.\\
\end{tabular}

\vfill
{\large \today}
\end{center}
\end{titlepage}

\tableofcontents
\newpage

%====================================================
\section*{Introduction générale}
\addcontentsline{toc}{section}{Introduction générale}

La chimie de première spécialité poursuit deux objectifs complémentaires. D'une part, elle cherche à décrire quantitativement un système chimique : combien de matière contient-il ? Comment cette quantité évolue-t-elle au cours d'une transformation ? Comment exploiter une transformation pour déterminer une concentration inconnue ? D'autre part, elle relie les propriétés macroscopiques des espèces chimiques à leur structure microscopique : géométrie, polarité, interactions, groupes caractéristiques et liaisons chimiques.

\begin{retenirbox}
La démarche générale en chimie consiste souvent à passer d'une description expérimentale à un modèle :
\[
\text{observations} \longrightarrow \text{espèces chimiques} \longrightarrow \text{équation de réaction} \longrightarrow \text{bilan quantitatif}.
\]
Ce passage demande de maîtriser les quantités de matière, les concentrations, les équations ajustées, les tableaux d'avancement et les raisonnements de proportionnalité.
\end{retenirbox}

%====================================================
\section{Suivi de l'évolution d'un système, siège d'une transformation}

\subsection{Quantité de matière, masse molaire et volume molaire}

\begin{definitionbox}
La \textbf{quantité de matière} d'un échantillon est notée $n$ et s'exprime en mole, de symbole \si{mol}. Une mole contient un nombre d'entités égal à la constante d'Avogadro :
\[
\Na \simeq \SI{6.02e23}{mol^{-1}}.
\]
Si un échantillon contient $N$ entités, alors :
\[
n=\frac{N}{\Na}.
\]
\end{definitionbox}

\begin{definitionbox}
La \textbf{masse molaire} $M$ d'une espèce chimique est la masse d'une mole de cette espèce. Elle s'exprime généralement en \si{g.mol^{-1}}. Pour un échantillon de masse $m$ :
\[
n=\frac{m}{M}, \qquad m=nM.
\]
\end{definitionbox}

\begin{methodbox}
Pour déterminer la masse molaire d'une molécule, on additionne les masses molaires atomiques des atomes qui la composent.

Par exemple, pour l'éthanol \ce{C2H6O} :
\[
M(\ce{C2H6O})=2M(\ce C)+6M(\ce H)+M(\ce O).
\]
Avec $M(\ce C)=\SI{12.0}{g.mol^{-1}}$, $M(\ce H)=\SI{1.0}{g.mol^{-1}}$ et $M(\ce O)=\SI{16.0}{g.mol^{-1}}$, on obtient :
\[
M(\ce{C2H6O})=2\times12.0+6\times1.0+16.0=\SI{46.0}{g.mol^{-1}}.
\]
\end{methodbox}

\begin{definitionbox}
Le \textbf{volume molaire} $V_m$ d'un gaz est le volume occupé par une mole de ce gaz dans des conditions données de température et de pression. Pour un gaz :
\[
n=\frac{V}{V_m}.
\]
Dans les conditions usuelles proches de \SI{20}{\degreeCelsius} et \SI{1}{bar}, on utilise souvent $V_m\simeq\SI{24.0}{L.mol^{-1}}$.
\end{definitionbox}

\begin{erreurbox}
Le volume molaire d'un gaz n'est pas une constante universelle indépendante des conditions expérimentales : il dépend de la température et de la pression. Dans un exercice, on utilise la valeur donnée dans l'énoncé.
\end{erreurbox}

\begin{exemplebox}
On dispose de \SI{5.80}{g} de butane \ce{C4H10}. Déterminer la quantité de matière correspondante.

\textbf{Solution.} La masse molaire vaut :
\[
M(\ce{C4H10})=4\times12.0+10\times1.0=\SI{58.0}{g.mol^{-1}}.
\]
La quantité de matière est donc :
\[
n=\frac{m}{M}=\frac{5.80}{58.0}=\SI{0.100}{mol}.
\]
L'échantillon contient \SI{0.100}{mol} de butane.
\end{exemplebox}

\subsection{Solutions et concentrations}

\begin{definitionbox}
Une \textbf{solution} est un mélange homogène composé d'un solvant et d'un ou plusieurs solutés. La \textbf{concentration en quantité de matière}, ou concentration molaire, d'un soluté est :
\[
c=\frac{n}{V},
\]
où $n$ est la quantité de matière de soluté et $V$ le volume de solution. Elle s'exprime en \si{mol.L^{-1}} si $V$ est exprimé en litres.
\end{definitionbox}

\begin{propositionbox}
Si la concentration en masse $C_m$ d'un soluté est connue, alors :
\[
C_m=\frac{m}{V}, \qquad c=\frac{C_m}{M}.
\]
\end{propositionbox}

\begin{demobox}
On sait que $m=nM$ et que $C_m=m/V$. En remplaçant $m$ par $nM$ :
\[
C_m=\frac{nM}{V}=M\frac{n}{V}=Mc.
\]
Donc :
\[
c=\frac{C_m}{M}.
\]
Cette relation illustre le lien entre une grandeur macroscopique mesurable, la masse dissoute, et une grandeur chimique fondamentale, la quantité de matière.
\end{demobox}

\begin{methodbox}
Pour résoudre un exercice de solution :
\begin{enumerate}
\item Identifier si l'on donne une masse, une concentration en masse ou une concentration molaire.
\item Convertir le volume en litres si la concentration est en \si{mol.L^{-1}}.
\item Utiliser $n=cV$, $m=C_mV$ ou $n=m/M$.
\item Vérifier l'unité du résultat.
\end{enumerate}
\end{methodbox}

\begin{exemplebox}
Une solution de sulfate de cuivre contient des ions \ce{Cu^{2+}} à la concentration $c=\SI{2.0e-2}{mol.L^{-1}}$. On prélève $V=\SI{50.0}{mL}$ de cette solution. Calculer la quantité de matière d'ions cuivre(II) prélevée.

\textbf{Solution.} On convertit le volume :
\[
V=\SI{50.0}{mL}=\SI{5.00e-2}{L}.
\]
Alors :
\[
n=cV=2.0\times10^{-2}\times5.00\times10^{-2}=\SI{1.0e-3}{mol}.
\]
\end{exemplebox}

\subsection{Absorbance, spectre d'absorption et loi de Beer-Lambert}

\begin{definitionbox}
L'\textbf{absorbance} $A$ d'une solution colorée mesure sa capacité à absorber une radiation lumineuse donnée. Elle dépend de la longueur d'onde utilisée, de la concentration de l'espèce colorée, de l'épaisseur traversée et de la nature de l'espèce.
\end{definitionbox}

\begin{theorembox}
Dans un domaine de validité expérimental, la loi de Beer-Lambert s'écrit :
\[
A=k c.
\]
Pour une longueur d'onde et une cuve données, l'absorbance est proportionnelle à la concentration de l'espèce colorée.
\end{theorembox}

\begin{demobox}
Au niveau de première, la loi de Beer-Lambert est admise comme un modèle expérimental. Sa justification qualitative est la suivante : plus une solution contient d'entités absorbantes, plus la probabilité qu'un photon soit absorbé est grande. Si l'on double la concentration et que l'on reste dans le domaine de validité, on double approximativement le nombre d'entités rencontrées par la lumière, donc l'absorbance.
\end{demobox}

\begin{methodbox}
Pour déterminer une concentration par étalonnage :
\begin{enumerate}
\item Préparer plusieurs solutions de concentrations connues.
\item Mesurer leur absorbance à une longueur d'onde adaptée, souvent proche du maximum d'absorption.
\item Tracer la droite d'étalonnage $A=f(c)$.
\item Mesurer l'absorbance de la solution inconnue.
\item Lire ou calculer sa concentration à partir de la droite.
\end{enumerate}
\end{methodbox}

\begin{erreurbox}
La couleur perçue d'une solution n'est pas la couleur absorbée. Une solution bleue absorbe plutôt les radiations complémentaires, notamment dans l'orange-rouge selon le cas. Il faut raisonner à partir du spectre d'absorption et non seulement à partir de l'apparence.
\end{erreurbox}

\subsection{Modéliser une transformation chimique}

\begin{definitionbox}
Une transformation chimique est modélisée par une \textbf{réaction chimique}, représentée par une équation ajustée. Les nombres placés devant les formules chimiques sont les \textbf{coefficients stœchiométriques}.
\end{definitionbox}

\begin{exemplebox}
La combustion complète du méthane est modélisée par :
\[
\ce{CH4 + 2 O2 -> CO2 + 2 H2O}.
\]
Cette équation signifie qu'une mole de méthane réagit avec deux moles de dioxygène pour former une mole de dioxyde de carbone et deux moles d'eau, si les réactifs sont introduits dans les proportions stœchiométriques.
\end{exemplebox}

\subsection{Oxydo-réduction}

\begin{definitionbox}
Un \textbf{oxydant} est une espèce capable de capter un ou plusieurs électrons. Un \textbf{réducteur} est une espèce capable de céder un ou plusieurs électrons.

Un couple oxydant-réducteur s'écrit :
\[
\text{Ox}/\text{Red}.
\]
La demi-équation électronique associée est de la forme :
\[
\text{Ox}+ne^- = \text{Red}.
\]
\end{definitionbox}

\begin{retenirbox}
Lors d'une réaction d'oxydo-réduction, il y a transfert d'électrons du réducteur vers l'oxydant.
\[
\text{réducteur donneur d'électrons} \quad \longrightarrow \quad \text{oxydant accepteur d'électrons}.
\]
\end{retenirbox}

\begin{methodbox}
Pour établir l'équation d'une réaction d'oxydo-réduction à partir de deux couples :
\begin{enumerate}
\item Écrire la demi-équation de réduction de l'oxydant.
\item Écrire la demi-équation d'oxydation du réducteur.
\item Multiplier les demi-équations pour échanger le même nombre d'électrons.
\item Additionner les demi-équations et simplifier les électrons.
\item Vérifier la conservation des éléments et des charges.
\end{enumerate}
\end{methodbox}

\begin{exemplebox}
On fait réagir des ions cuivre(II) \ce{Cu^{2+}} avec du zinc métallique \ce{Zn}. Les couples sont \ce{Cu^{2+}/Cu} et \ce{Zn^{2+}/Zn}.

\textbf{Demi-équations :}
\[
\ce{Cu^{2+} + 2e^- -> Cu}
\]
\[
\ce{Zn -> Zn^{2+} + 2e^-}
\]
En additionnant :
\[
\ce{Cu^{2+} + Zn -> Cu + Zn^{2+}}.
\]
Les électrons n'apparaissent pas dans l'équation globale, car ils sont transférés entre les espèces.
\end{exemplebox}

\subsection{Avancement et tableau d'avancement}

\begin{definitionbox}
L'\textbf{avancement} $x$ d'une réaction est une grandeur exprimée en mole qui permet de suivre l'évolution des quantités de matière au cours d'une transformation. Il traduit le nombre de fois où la réaction, telle qu'elle est écrite, s'est produite à l'échelle macroscopique.
\end{definitionbox}

\begin{methodbox}
Pour une réaction générale :
\[
aA+bB\rightarrow cC+dD,
\]
si l'avancement vaut $x$, les variations de quantité de matière sont :
\[
\Delta n(A)=-ax,\quad \Delta n(B)=-bx,\quad \Delta n(C)=+cx,\quad \Delta n(D)=+dx.
\]
\end{methodbox}

\begin{theorembox}
L'avancement maximal $\xmax$ est la plus grande valeur possible de $x$ compatible avec des quantités de matière finales positives ou nulles. Il est atteint lorsque le réactif limitant est totalement consommé.
\end{theorembox}

\begin{demobox}
Pour que le tableau d'avancement ait un sens physique, aucune quantité de matière finale ne peut être négative. Si $n_i(A)-ax\geq0$, alors $x\leq n_i(A)/a$. De même, $x\leq n_i(B)/b$. Donc :
\[
\xmax=\min\left(\frac{n_i(A)}{a},\frac{n_i(B)}{b}\right).
\]
Le plus petit quotient correspond au réactif qui impose l'arrêt de la transformation totale : c'est le réactif limitant.
\end{demobox}

\begin{exemplebox}
On mélange $n_i(\ce{H2})=\SI{0.30}{mol}$ et $n_i(\ce{O2})=\SI{0.10}{mol}$. La réaction est :
\[
\ce{2 H2 + O2 -> 2 H2O}.
\]
Le tableau d'avancement est :
\[
\begin{array}{c|ccc}
 & \ce{H2} & \ce{O2} & \ce{H2O} \\
\hline
\text{initial} & 0.30 & 0.10 & 0 \\
\text{variation} & -2x & -x & +2x \\
\text{final} & 0.30-2x & 0.10-x & 2x
\end{array}
\]
Les contraintes sont :
\[
0.30-2x\geq0 \Rightarrow x\leq0.15,\qquad 0.10-x\geq0 \Rightarrow x\leq0.10.
\]
Donc $\xmax=\SI{0.10}{mol}$ et \ce{O2} est limitant.
\end{exemplebox}

\subsection{Transformation totale, non totale et mélange stœchiométrique}

\begin{definitionbox}
Une transformation est dite \textbf{totale} si son avancement final est égal à l'avancement maximal :
\[
\xf=\xmax.
\]
Elle est dite \textbf{non totale} si :
\[
\xf<\xmax.
\]
Au niveau de première, on compare simplement ces deux valeurs, sans introduire la notion d'équilibre chimique.
\end{definitionbox}

\begin{definitionbox}
Un mélange initial est \textbf{stœchiométrique} lorsque les réactifs sont introduits dans les proportions des coefficients de l'équation. Pour :
\[
aA+bB\rightarrow cC+dD,
\]
un mélange est stœchiométrique si :
\[
\frac{n_i(A)}{a}=\frac{n_i(B)}{b}.
\]
\end{definitionbox}

\subsection{Titrage avec suivi colorimétrique}

\begin{definitionbox}
Un \textbf{titrage} est une méthode expérimentale permettant de déterminer la quantité de matière ou la concentration d'une espèce chimique en la faisant réagir avec une espèce de concentration connue, appelée solution titrante.
\end{definitionbox}

\begin{definitionbox}
L'\textbf{équivalence} d'un titrage est l'instant où les réactifs titré et titrant ont été introduits dans les proportions stœchiométriques. À l'équivalence, il y a changement de réactif limitant.
\end{definitionbox}

\begin{methodbox}
Pour exploiter un titrage :
\begin{enumerate}
\item Écrire l'équation de la réaction support du titrage.
\item Identifier les coefficients stœchiométriques.
\item À l'équivalence, écrire la relation de proportionnalité :
\[
\frac{n(\text{espèce titrée})}{a}=\frac{n(\text{espèce titrante versée})}{b}.
\]
\item Remplacer $n$ par $cV$ pour les solutions.
\item Isoler la concentration inconnue.
\end{enumerate}
\end{methodbox}

\begin{exemplebox}
On titre $V_A=\SI{10.0}{mL}$ d'une solution contenant des ions fer(II) \ce{Fe^{2+}} par une solution de permanganate \ce{MnO4^-} de concentration $c_B=\SI{2.00e-2}{mol.L^{-1}}$. Le volume équivalent est $V_E=\SI{12.0}{mL}$. La réaction en milieu acide est :
\[
\ce{MnO4^- + 5 Fe^{2+} + 8H+ -> Mn^{2+} + 5Fe^{3+} + 4H2O}.
\]
À l'équivalence :
\[
\frac{n(\ce{Fe^{2+}})}{5}=n(\ce{MnO4^-}).
\]
Donc :
\[
c_A V_A =5c_B V_E.
\]
Avec les volumes en litres :
\[
c_A=\frac{5c_BV_E}{V_A}=\frac{5\times2.00\times10^{-2}\times12.0\times10^{-3}}{10.0\times10^{-3}}
=\SI{0.120}{mol.L^{-1}}.
\]
\end{exemplebox}

\subsection{Exercices progressifs - Partie 1}

\begin{exercicebox}[title=Exercice 1 - Quantité de matière]
On dispose de \SI{9.00}{g} de glucose \ce{C6H12O6}.
\begin{enumerate}
\item Calculer la masse molaire du glucose.
\item Déterminer la quantité de matière correspondante.
\item Calculer le nombre de molécules de glucose.
\end{enumerate}
Donnée : $\Na=\SI{6.02e23}{mol^{-1}}$.
\end{exercicebox}

\begin{exercicebox}[title=Exercice 2 - Solution et concentration]
On dissout \SI{2.34}{g} de chlorure de sodium \ce{NaCl} dans une fiole jaugée de volume \SI{250.0}{mL}.
\begin{enumerate}
\item Calculer la masse molaire de \ce{NaCl}.
\item Déterminer la quantité de matière dissoute.
\item Calculer la concentration molaire de la solution.
\end{enumerate}
Données : $M(\ce{Na})=\SI{23.0}{g.mol^{-1}}$, $M(\ce{Cl})=\SI{35.5}{g.mol^{-1}}$.
\end{exercicebox}

\begin{exercicebox}[title=Exercice 3 - Tableau d'avancement]
On fait réagir \SI{0.200}{mol} d'aluminium avec \SI{0.300}{mol} d'ions cuivre(II). La réaction est :
\[
\ce{2Al + 3Cu^{2+} -> 2Al^{3+} + 3Cu}.
\]
\begin{enumerate}
\item Construire le tableau d'avancement.
\item Déterminer l'avancement maximal et le réactif limitant.
\item Déterminer la composition finale si la transformation est totale.
\end{enumerate}
\end{exercicebox}

\begin{exercicebox}[title=Exercice 4 - Titrage d'une eau oxygénée]
On titre $V_A=\SI{10.0}{mL}$ d'une solution d'eau oxygénée contenant \ce{H2O2} par une solution de permanganate de potassium acidifiée de concentration $c_B=\SI{2.00e-2}{mol.L^{-1}}$. L'équivalence est atteinte pour $V_E=\SI{8.60}{mL}$. L'équation support est :
\[
\ce{2MnO4^- + 5H2O2 + 6H+ -> 2Mn^{2+} + 5O2 + 8H2O}.
\]
Déterminer la concentration molaire de \ce{H2O2} dans la solution titrée.
\end{exercicebox}

\newpage

%====================================================
\section{De la structure des entités aux propriétés physiques de la matière}

\subsection{Schémas de Lewis}

\begin{definitionbox}
Un \textbf{schéma de Lewis} représente les atomes d'une entité chimique, les liaisons covalentes et les doublets non liants portés par les atomes. Il permet de visualiser la répartition des électrons de valence.
\end{definitionbox}

\begin{retenirbox}
Pour les atomes usuels en première :
\begin{center}
\begin{tabular}{c|c|c}
Atome & Électrons de valence usuels & Nombre de liaisons courant \\
\hline
\ce H & 1 & 1 \\
\ce C & 4 & 4 \\
\ce N & 5 & 3, avec un doublet non liant \\
\ce O & 6 & 2, avec deux doublets non liants \\
\ce Cl & 7 & 1, avec trois doublets non liants
\end{tabular}
\end{center}
\end{retenirbox}

\begin{methodbox}
Pour établir un schéma de Lewis simple :
\begin{enumerate}
\item Compter les électrons de valence de tous les atomes, en tenant compte de la charge de l'ion.
\item Choisir un squelette de l'entité, généralement l'atome le moins électronégatif au centre, sauf l'hydrogène qui est toujours périphérique.
\item Placer les liaisons simples.
\item Compléter les octets des atomes périphériques.
\item Placer les électrons restants sur l'atome central.
\item Si l'atome central n'a pas l'octet, former des liaisons multiples si nécessaire.
\end{enumerate}
\end{methodbox}

\begin{definitionbox}
Une \textbf{lacune électronique} est un manque de doublet autour d'un atome. Une entité qui possède une lacune électronique peut accepter un doublet d'électrons.
\end{definitionbox}

\subsection{Géométrie des entités}

\begin{propositionbox}
La géométrie d'une entité dépend de la répartition des doublets d'électrons autour de l'atome central. Les doublets, liants ou non liants, se repoussent et se placent de façon à minimiser les répulsions.
\end{propositionbox}

\begin{center}
\begin{tabularx}{\textwidth}{c|c|X}
\toprule
Nombre de domaines autour du centre & Exemple & Géométrie approximative \\
\midrule
2 & \ce{CO2} & linéaire \\
3 & \ce{CH2O} & triangulaire plane autour du carbone \\
4 & \ce{CH4} & tétraédrique \\
4 dont 1 doublet non liant & \ce{NH3} & pyramidale \\
4 dont 2 doublets non liants & \ce{H2O} & coudée \\
\bottomrule
\end{tabularx}
\end{center}

\begin{erreurbox}
Un schéma de Lewis est plan sur la feuille, mais la molécule réelle peut être tridimensionnelle. Par exemple, \ce{CH4} n'est pas une croix plane : sa géométrie est tétraédrique.
\end{erreurbox}

\subsection{Électronégativité, polarisation et polarité}

\begin{definitionbox}
L'\textbf{électronégativité} traduit la tendance d'un atome engagé dans une liaison covalente à attirer vers lui les électrons de la liaison.
\end{definitionbox}

\begin{retenirbox}
Dans le tableau périodique, l'électronégativité augmente globalement de gauche à droite dans une période et diminue lorsque l'on descend dans une famille. Le fluor est l'élément le plus électronégatif.
\end{retenirbox}

\begin{definitionbox}
Une liaison covalente est \textbf{polarisée} si les deux atomes liés n'ont pas la même électronégativité. L'atome le plus électronégatif porte une charge partielle négative $\delta^-$ et l'autre une charge partielle positive $\delta^+$.
\end{definitionbox}

\begin{definitionbox}
Une molécule est \textbf{polaire} si les polarités de ses liaisons ne se compensent pas dans l'espace. Elle est \textbf{apolaire} si les polarités se compensent ou si ses liaisons ne sont pas polarisées.
\end{definitionbox}

\begin{methodbox}
Pour déterminer si une molécule est polaire :
\begin{enumerate}
\item Étudier la polarité de chaque liaison.
\item Déterminer la géométrie de la molécule.
\item Examiner si les effets de polarisation se compensent ou non.
\end{enumerate}
\end{methodbox}

\begin{exemplebox}
La molécule \ce{CO2} possède deux liaisons \ce{C=O} polarisées, mais elle est linéaire :
\[
\ce{O=C=O}.
\]
Les deux polarités sont opposées et se compensent : \ce{CO2} est apolaire.

La molécule \ce{H2O} possède deux liaisons \ce{O-H} polarisées, et sa géométrie est coudée. Les polarités ne se compensent pas : \ce{H2O} est polaire.
\end{exemplebox}

\subsection{Interactions et cohésion de la matière}

\begin{definitionbox}
La \textbf{cohésion} d'un solide ou d'un liquide résulte des interactions entre les entités microscopiques qui le constituent : ions, molécules polaires, molécules apolaires, liaisons hydrogène.
\end{definitionbox}

\begin{retenirbox}
On distingue notamment :
\begin{itemize}
\item les interactions électrostatiques entre ions, fortes ;
\item les interactions entre molécules polaires ;
\item les interactions entre molécules apolaires, plus faibles mais nombreuses ;
\item les ponts hydrogène, qui sont des interactions particulièrement importantes lorsqu'un hydrogène lié à \ce O, \ce N ou \ce F interagit avec un doublet non liant d'un autre atome très électronégatif.
\end{itemize}
\end{retenirbox}

\begin{exemplebox}
L'eau possède une température d'ébullition élevée par rapport à des molécules de taille comparable. Cela s'explique par la présence de nombreux ponts hydrogène entre les molécules d'eau.
\end{exemplebox}

\subsection{Dissolution des solides ioniques dans l'eau}

\begin{definitionbox}
La \textbf{dissolution} d'un solide ionique dans l'eau correspond à la dissociation du solide en ions, puis à la solvatation de ces ions par les molécules d'eau.
\end{definitionbox}

\begin{propositionbox}
L'eau dissout de nombreux solides ioniques car elle est polaire. Le pôle négatif de la molécule d'eau entoure les cations, tandis que les pôles positifs entourent les anions.
\end{propositionbox}

\begin{exemplebox}
La dissolution du chlorure de sodium dans l'eau s'écrit :
\[
\ce{NaCl(s) -> Na+(aq) + Cl-(aq)}.
\]
Si l'on dissout \SI{0.100}{mol} de \ce{NaCl} dans \SI{1.00}{L} d'eau, alors :
\[
[\ce{Na+}]=\SI{0.100}{mol.L^{-1}},\qquad [\ce{Cl-}]=\SI{0.100}{mol.L^{-1}}.
\]
\end{exemplebox}

\begin{exemplebox}
Pour le chlorure de calcium :
\[
\ce{CaCl2(s) -> Ca^{2+}(aq) + 2Cl-(aq)}.
\]
Si la concentration apportée en \ce{CaCl2} vaut $c=\SI{0.050}{mol.L^{-1}}$, alors :
\[
[\ce{Ca^{2+}}]=c=\SI{0.050}{mol.L^{-1}},
\]
mais :
\[
[\ce{Cl-}]=2c=\SI{0.100}{mol.L^{-1}}.
\]
\end{exemplebox}

\subsection{Solubilité, miscibilité et extraction liquide-liquide}

\begin{definitionbox}
La \textbf{solubilité} d'une espèce dans un solvant correspond à la quantité maximale que l'on peut dissoudre dans un volume donné de solvant, à température donnée.
\end{definitionbox}

\begin{definitionbox}
Deux liquides sont \textbf{miscibles} s'ils forment un mélange homogène en toutes proportions. Ils sont \textbf{non miscibles} s'ils forment deux phases distinctes.
\end{definitionbox}

\begin{retenirbox}
Principe qualitatif : une espèce se dissout mieux dans un solvant avec lequel elle peut établir des interactions favorables. On résume souvent par :
\[
\text{le semblable dissout le semblable}.
\]
Une espèce polaire est souvent plus soluble dans un solvant polaire ; une espèce apolaire est souvent plus soluble dans un solvant apolaire.
\end{retenirbox}

\begin{methodbox}
Pour choisir un solvant d'extraction :
\begin{enumerate}
\item Le solvant doit bien dissoudre l'espèce à extraire.
\item Il doit être non miscible avec le solvant initial.
\item Il ne doit pas réagir avec l'espèce à extraire.
\item Il doit si possible être peu toxique, peu inflammable et facile à éliminer.
\end{enumerate}
\end{methodbox}

\subsection{Hydrophilie, lipophilie et amphiphilie}

\begin{definitionbox}
Une espèce \textbf{hydrophile} interagit favorablement avec l'eau. Une espèce \textbf{lipophile} interagit favorablement avec les corps gras. Une espèce \textbf{amphiphile} possède une partie hydrophile et une partie lipophile.
\end{definitionbox}

\begin{exemplebox}
Un ion carboxylate à longue chaîne, comme un savon, possède une tête ionique hydrophile et une longue chaîne carbonée lipophile. Il peut donc entourer des graisses et permettre leur dispersion dans l'eau.
\end{exemplebox}

\begin{retenirbox}
Les tensioactifs diminuent la tension superficielle de l'eau et favorisent la dispersion de substances grasses sous forme de micelles. Ils sont utilisés dans les savons, détergents, lessives, shampoings et certains procédés industriels.
\end{retenirbox}

\subsection{Exercices progressifs - Partie 2}

\begin{exercicebox}[title=Exercice 5 - Schémas de Lewis]
Établir les schémas de Lewis des entités suivantes : \ce{H2O}, \ce{NH3}, \ce{CO2}, \ce{CH4}, \ce{Cl-}. Indiquer pour chacune le nombre de doublets liants et non liants autour de l'atome central lorsqu'il existe.
\end{exercicebox}

\begin{exercicebox}[title=Exercice 6 - Polarité]
On considère les molécules \ce{H2O}, \ce{CO2}, \ce{NH3} et \ce{CH4}.
\begin{enumerate}
\item Indiquer la géométrie de chacune.
\item Dire si les liaisons sont polarisées.
\item Conclure sur le caractère polaire ou apolaire de chaque molécule.
\end{enumerate}
\end{exercicebox}

\begin{exercicebox}[title=Exercice 7 - Dissolution ionique]
On dissout \SI{0.0200}{mol} de sulfate de sodium \ce{Na2SO4} dans de l'eau pour obtenir \SI{200.0}{mL} de solution.
\begin{enumerate}
\item Écrire l'équation de dissolution.
\item Calculer la concentration en ions \ce{Na+}.
\item Calculer la concentration en ions \ce{SO4^{2-}}.
\end{enumerate}
\end{exercicebox}

\begin{exercicebox}[title=Exercice 8 - Extraction]
Une espèce organique $X$ est soluble dans l'eau à raison de \SI{2.0}{g.L^{-1}} et dans le cyclohexane à raison de \SI{35}{g.L^{-1}}. L'eau et le cyclohexane ne sont pas miscibles.
\begin{enumerate}
\item Quel solvant semble le plus adapté pour extraire $X$ d'une phase aqueuse ?
\item Expliquer le rôle de l'ampoule à décanter.
\item Quelles informations faut-il connaître pour savoir quelle phase est au-dessus ?
\end{enumerate}
\end{exercicebox}

\newpage

%====================================================
\section{Propriétés physico-chimiques, synthèses et combustions d'espèces organiques}

\subsection{Formules et chaînes carbonées}

\begin{definitionbox}
Une molécule organique contient généralement des atomes de carbone et d'hydrogène, éventuellement associés à d'autres atomes comme \ce O, \ce N, \ce Cl. La \textbf{formule brute} donne la nature et le nombre des atomes. La \textbf{formule semi-développée} montre l'enchaînement des atomes sans représenter toutes les liaisons \ce{C-H} séparément.
\end{definitionbox}

\begin{definitionbox}
Un squelette carboné est dit \textbf{saturé} lorsque les atomes de carbone sont reliés uniquement par des liaisons simples. Les alcanes sont des hydrocarbures saturés de formule brute générale \ce{C_nH_{2n+2}}.
\end{definitionbox}

\subsection{Groupes caractéristiques et familles fonctionnelles}

\begin{retenirbox}
\begin{center}
\begin{tabularx}{\textwidth}{c|c|X}
\toprule
Famille & Groupe caractéristique & Exemple \\
\midrule
Alcool & hydroxyle \ce{-OH} porté par un carbone saturé & éthanol \ce{CH3-CH2-OH} \\
Aldéhyde & carbonyle terminal \ce{-CHO} & éthanal \ce{CH3-CHO} \\
Cétone & carbonyle interne $>\ce{C=O}$ & propanone \ce{CH3-CO-CH3} \\
Acide carboxylique & carboxyle \ce{-COOH} & acide éthanoïque \ce{CH3-COOH} \\
\bottomrule
\end{tabularx}
\end{center}
\end{retenirbox}

\begin{methodbox}
Pour nommer une molécule organique simple avec un seul groupe caractéristique :
\begin{enumerate}
\item Repérer la chaîne carbonée principale.
\item Compter le nombre de carbones : méth-, éth-, prop-, but-, pent-, etc.
\item Identifier la famille fonctionnelle.
\item Placer si nécessaire le numéro du carbone portant le groupe caractéristique.
\item Ajouter le suffixe adapté : -ol, -al, -one, acide ... -oïque.
\end{enumerate}
\end{methodbox}

\begin{exemplebox}
La formule semi-développée \ce{CH3-CH2-CH2-OH} possède trois atomes de carbone et un groupe alcool. Son nom est propan-1-ol.

La formule \ce{CH3-CO-CH3} possède trois carbones et un groupe cétone au carbone 2. Son nom est propanone.
\end{exemplebox}

\subsection{Spectroscopie infrarouge}

\begin{definitionbox}
La spectroscopie infrarouge permet d'identifier certains groupes caractéristiques grâce à l'absorption de radiations infrarouges correspondant à des vibrations de liaisons.
\end{definitionbox}

\begin{retenirbox}
Quelques repères usuels :
\begin{center}
\begin{tabular}{c|c|c}
Liaison / groupe & Nombre d'onde approximatif & Aspect \\
\hline
\ce{O-H} alcool & \SIrange{3200}{3600}{cm^{-1}} & bande large \\
\ce{C=O} & autour de \SI{1700}{cm^{-1}} & bande forte \\
\ce{C-H} & autour de \SI{2900}{cm^{-1}} & bandes \\
\ce{O-H} acide carboxylique & \SIrange{2500}{3300}{cm^{-1}} & très large \\
\end{tabular}
\end{center}
\end{retenirbox}

\subsection{Synthèse organique et rendement}

\begin{definitionbox}
Une synthèse organique suit souvent plusieurs étapes : transformation chimique des réactifs, isolement du produit, purification, puis analyse ou identification.
\end{definitionbox}

\begin{definitionbox}
Le \textbf{rendement} d'une synthèse est :
\[
\eta=\frac{n_{\text{obtenu}}}{n_{\text{maximal théorique}}}\times100.
\]
On peut aussi utiliser les masses si le produit comparé est le même :
\[
\eta=\frac{m_{\text{obtenu}}}{m_{\text{maximal théorique}}}\times100.
\]
\end{definitionbox}

\begin{methodbox}
Pour calculer un rendement :
\begin{enumerate}
\item Écrire l'équation de la réaction.
\item Déterminer le réactif limitant.
\item Calculer la quantité maximale théorique de produit.
\item Comparer à la quantité réellement obtenue.
\item Donner le rendement en pourcentage.
\end{enumerate}
\end{methodbox}

\begin{erreurbox}
Un rendement ne peut pas dépasser \SI{100}{\percent} dans un résultat cohérent. Un rendement supérieur à \SI{100}{\percent} signale souvent une erreur de calcul, un produit humide, impur, ou une masse expérimentale mal interprétée.
\end{erreurbox}

\subsection{Combustions d'espèces organiques}

\begin{definitionbox}
Une combustion complète est une réaction avec le dioxygène qui transforme un composé organique contenant du carbone et de l'hydrogène en dioxyde de carbone et en eau.
\end{definitionbox}

\begin{methodbox}
Pour ajuster la combustion complète d'un alcane \ce{C_nH_{2n+2}} :
\begin{enumerate}
\item Placer \ce{C_nH_{2n+2}} et \ce{O2} à gauche.
\item Placer $n\ce{CO2}$ à droite pour conserver le carbone.
\item Placer $(n+1)\ce{H2O}$ à droite pour conserver l'hydrogène.
\item Ajuster le dioxygène.
\end{enumerate}
\end{methodbox}

\begin{exemplebox}
Combustion complète du propane :
\[
\ce{C3H8 + 5O2 -> 3CO2 + 4H2O}.
\]
On vérifie : 3 atomes de carbone de chaque côté, 8 hydrogènes de chaque côté, 10 oxygènes de chaque côté.
\end{exemplebox}

\subsection{Énergie de liaison et énergie molaire de réaction}

\begin{definitionbox}
L'\textbf{énergie de liaison} est l'énergie qu'il faut fournir pour rompre une mole d'une liaison donnée en phase gazeuse. Elle s'exprime souvent en \si{kJ.mol^{-1}}.
\end{definitionbox}

\begin{theorembox}
Une estimation de l'énergie molaire de réaction en phase gazeuse est :
\[
\Delta_r E \approx \sum E_{\text{liaisons rompues}}-\sum E_{\text{liaisons formées}}.
\]
Si $\Delta_rE<0$, la transformation libère de l'énergie : elle est exothermique.
\end{theorembox}

\begin{demobox}
Rompre des liaisons demande de fournir de l'énergie : contribution positive. Former des liaisons libère de l'énergie : contribution négative. Le bilan énergétique est donc l'énergie nécessaire aux ruptures moins l'énergie libérée lors des formations.
\end{demobox}

\begin{definitionbox}
Le \textbf{pouvoir calorifique massique} d'un combustible est l'énergie libérée par la combustion complète d'un kilogramme de combustible. Il s'exprime en \si{J.kg^{-1}} ou \si{MJ.kg^{-1}}.
\end{definitionbox}

\subsection{Combustions et enjeux de société}

Les combustions sont utilisées pour le chauffage, le transport, la production d'électricité ou certains procédés industriels. Elles présentent cependant des risques : incendie, explosion, intoxication au monoxyde de carbone lors d'une combustion incomplète, émission de dioxyde de carbone et de polluants.

\begin{retenirbox}
Les axes actuels d'étude incluent :
\begin{itemize}
\item la réduction des émissions de gaz à effet de serre ;
\item le développement de carburants moins carbonés ;
\item l'amélioration des rendements énergétiques ;
\item la valorisation de la biomasse et du biogaz ;
\item la sobriété énergétique et la limitation des combustions inutiles.
\end{itemize}
\end{retenirbox}

\subsection{Exercices progressifs - Partie 3}

\begin{exercicebox}[title=Exercice 9 - Groupes caractéristiques]
Identifier la famille chimique des molécules suivantes :
\[
\ce{CH3-CH2-OH},\qquad \ce{CH3-CHO},\qquad \ce{CH3-CO-CH3},\qquad \ce{CH3-COOH}.
\]
Donner le nom usuel ou systématique simple de chacune.
\end{exercicebox}

\begin{exercicebox}[title=Exercice 10 - Rendement d'une synthèse]
Une synthèse permet théoriquement d'obtenir \SI{8.50}{g} d'un ester. Après purification, on recueille \SI{6.20}{g} de produit pur.
\begin{enumerate}
\item Calculer le rendement.
\item Proposer deux raisons expérimentales pouvant expliquer un rendement inférieur à \SI{100}{\percent}.
\end{enumerate}
\end{exercicebox}

\begin{exercicebox}[title=Exercice 11 - Combustion complète]
Écrire l'équation ajustée de la combustion complète :
\begin{enumerate}
\item du méthane \ce{CH4} ;
\item de l'éthanol \ce{C2H6O} ;
\item du butane \ce{C4H10}.
\end{enumerate}
\end{exercicebox}

\begin{exercicebox}[title=Exercice 12 - Énergie de combustion]
On modélise la combustion du méthane par :
\[
\ce{CH4 + 2O2 -> CO2 + 2H2O}.
\]
On donne les énergies de liaison :
\[
E(\ce{C-H})=\SI{413}{kJ.mol^{-1}},\quad E(\ce{O=O})=\SI{498}{kJ.mol^{-1}},
\]
\[
E(\ce{C=O})=\SI{799}{kJ.mol^{-1}},\quad E(\ce{O-H})=\SI{463}{kJ.mol^{-1}}.
\]
Estimer l'énergie molaire de réaction.
\end{exercicebox}

\newpage

%====================================================
\section{Grand devoir bilan final}

\begin{exercicebox}[title=Devoir bilan - Partie A : dosage et transformation chimique]
On souhaite déterminer la concentration en ions fer(II) \ce{Fe^{2+}} dans une solution $S$. On prélève $V_A=\SI{20.0}{mL}$ de cette solution et on la titre par une solution acidifiée de permanganate de potassium contenant des ions \ce{MnO4^-} de concentration $c_B=\SI{1.00e-2}{mol.L^{-1}}$. L'équivalence est observée pour $V_E=\SI{16.4}{mL}$.

L'équation support du titrage est :
\[
\ce{MnO4^- + 5Fe^{2+} + 8H+ -> Mn^{2+} + 5Fe^{3+} + 4H2O}.
\]
\begin{enumerate}
\item Identifier l'oxydant et le réducteur parmi \ce{MnO4^-} et \ce{Fe^{2+}}.
\item Écrire la relation entre les quantités de matière à l'équivalence.
\item Déterminer la concentration $c_A$ en ions \ce{Fe^{2+}} dans la solution $S$.
\item Expliquer le changement de réactif limitant à l'équivalence.
\end{enumerate}
\end{exercicebox}

\begin{exercicebox}[title=Devoir bilan - Partie B : structure et solubilité]
On considère l'éthanol \ce{CH3-CH2-OH}, le dioxyde de carbone \ce{CO2} et l'eau \ce{H2O}.
\begin{enumerate}
\item Identifier le groupe caractéristique de l'éthanol.
\item Expliquer pourquoi l'éthanol est miscible avec l'eau.
\item Comparer qualitativement la polarité de \ce{CO2} et \ce{H2O}.
\item Expliquer pourquoi l'eau dissout bien de nombreux solides ioniques.
\end{enumerate}
\end{exercicebox}

\begin{exercicebox}[title=Devoir bilan - Partie C : synthèse et combustion]
On réalise une synthèse qui peut théoriquement produire \SI{0.120}{mol} d'un ester de masse molaire $M=\SI{88.0}{g.mol^{-1}}$. On obtient expérimentalement \SI{7.80}{g} de produit pur.
\begin{enumerate}
\item Calculer la masse théorique maximale de produit.
\item Calculer le rendement de la synthèse.
\item L'ester contient seulement C, H et O. Lors de sa combustion complète, quels produits se forment ?
\item Citer un risque associé aux combustions et un enjeu environnemental.
\end{enumerate}
\end{exercicebox}

\newpage

%====================================================
\section{Corrections détaillées des exercices}

\subsection*{Correction exercice 1}
\begin{correctionbox}
La formule du glucose est \ce{C6H12O6}. Sa masse molaire vaut :
\[
M=6\times12.0+12\times1.0+6\times16.0=72.0+12.0+96.0=\SI{180.0}{g.mol^{-1}}.
\]
La quantité de matière est :
\[
n=\frac{m}{M}=\frac{9.00}{180.0}=\SI{5.00e-2}{mol}.
\]
Le nombre de molécules est :
\[
N=n\Na=5.00\times10^{-2}\times6.02\times10^{23}=\SI{3.01e22}{}.
\]
Il y a donc environ $\SI{3.01e22}{}$ molécules de glucose.
\end{correctionbox}

\subsection*{Correction exercice 2}
\begin{correctionbox}
La masse molaire du chlorure de sodium vaut :
\[
M(\ce{NaCl})=23.0+35.5=\SI{58.5}{g.mol^{-1}}.
\]
La quantité de matière dissoute est :
\[
n=\frac{2.34}{58.5}=\SI{4.00e-2}{mol}.
\]
Le volume vaut :
\[
V=\SI{250.0}{mL}=\SI{0.2500}{L}.
\]
La concentration est donc :
\[
c=\frac{n}{V}=\frac{4.00\times10^{-2}}{0.2500}=\SI{0.160}{mol.L^{-1}}.
\]
\end{correctionbox}

\subsection*{Correction exercice 3}
\begin{correctionbox}
Pour :
\[
\ce{2Al + 3Cu^{2+} -> 2Al^{3+} + 3Cu},
\]
le tableau d'avancement est :
\[
\begin{array}{c|cccc}
 & \ce{Al} & \ce{Cu^{2+}} & \ce{Al^{3+}} & \ce{Cu} \\
\hline
\text{initial} & 0.200 & 0.300 & 0 & 0 \\
\text{variation} & -2x & -3x & +2x & +3x \\
\text{final} & 0.200-2x & 0.300-3x & 2x & 3x
\end{array}
\]
Contraintes :
\[
0.200-2x\geq0\Rightarrow x\leq0.100,
\]
\[
0.300-3x\geq0\Rightarrow x\leq0.100.
\]
Donc $\xmax=\SI{0.100}{mol}$. Les deux réactifs sont totalement consommés : le mélange est stœchiométrique. À l'état final :
\[
n_f(\ce{Al})=0,
\quad n_f(\ce{Cu^{2+}})=0,
\quad n_f(\ce{Al^{3+}})=\SI{0.200}{mol},
\quad n_f(\ce{Cu})=\SI{0.300}{mol}.
\]
\end{correctionbox}

\subsection*{Correction exercice 4}
\begin{correctionbox}
À l'équivalence, les réactifs sont introduits dans les proportions stœchiométriques :
\[
\frac{n(\ce{H2O2})}{5}=\frac{n(\ce{MnO4^-})}{2}.
\]
Donc :
\[
n(\ce{H2O2})=\frac{5}{2}n(\ce{MnO4^-}).
\]
Or :
\[
n(\ce{MnO4^-})=c_BV_E=2.00\times10^{-2}\times8.60\times10^{-3}=\SI{1.72e-4}{mol}.
\]
Ainsi :
\[
n(\ce{H2O2})=\frac{5}{2}\times1.72\times10^{-4}=\SI{4.30e-4}{mol}.
\]
Comme $V_A=\SI{10.0}{mL}=\SI{1.00e-2}{L}$ :
\[
c_A=\frac{n}{V_A}=\frac{4.30\times10^{-4}}{1.00\times10^{-2}}=\SI{4.30e-2}{mol.L^{-1}}.
\]
\end{correctionbox}

\subsection*{Correction exercice 5}
\begin{correctionbox}
\ce{H2O} : l'oxygène est central, il forme deux liaisons \ce{O-H} et porte deux doublets non liants.

\ce{NH3} : l'azote est central, il forme trois liaisons \ce{N-H} et porte un doublet non liant.

\ce{CO2} : le carbone est central, il forme deux doubles liaisons \ce{C=O}. Chaque oxygène porte deux doublets non liants.

\ce{CH4} : le carbone est central et forme quatre liaisons \ce{C-H}. Il ne porte pas de doublet non liant.

\ce{Cl-} : l'ion chlorure possède quatre doublets autour du chlore ; il est monoatomique et porte une charge négative.
\end{correctionbox}

\subsection*{Correction exercice 6}
\begin{correctionbox}
\ce{H2O} est coudée. Les liaisons \ce{O-H} sont polarisées vers O. Les polarités ne se compensent pas : molécule polaire.

\ce{CO2} est linéaire. Les liaisons \ce{C=O} sont polarisées, mais leurs effets se compensent : molécule apolaire.

\ce{NH3} est pyramidale. Les liaisons \ce{N-H} sont polarisées vers N. Les polarités ne se compensent pas : molécule polaire.

\ce{CH4} est tétraédrique. Les liaisons \ce{C-H} sont peu polarisées et la géométrie est symétrique : molécule apolaire.
\end{correctionbox}

\subsection*{Correction exercice 7}
\begin{correctionbox}
L'équation de dissolution est :
\[
\ce{Na2SO4(s) -> 2Na+(aq) + SO4^{2-}(aq)}.
\]
Le volume vaut $V=\SI{200.0}{mL}=\SI{0.2000}{L}$. La concentration apportée en sulfate de sodium est :
\[
c=\frac{0.0200}{0.2000}=\SI{0.100}{mol.L^{-1}}.
\]
D'après l'équation :
\[
[\ce{SO4^{2-}}]=c=\SI{0.100}{mol.L^{-1}},
\]
\[
[\ce{Na+}]=2c=\SI{0.200}{mol.L^{-1}}.
\]
\end{correctionbox}

\subsection*{Correction exercice 8}
\begin{correctionbox}
L'espèce $X$ est beaucoup plus soluble dans le cyclohexane que dans l'eau : \SI{35}{g.L^{-1}} contre \SI{2.0}{g.L^{-1}}. Le cyclohexane est donc adapté pour extraire $X$ d'une phase aqueuse, à condition qu'il ne réagisse pas avec $X$.

L'ampoule à décanter permet de mélanger les deux phases, puis de les laisser se séparer selon leur non-miscibilité. L'espèce $X$ passe préférentiellement dans la phase où elle est la plus soluble.

Pour savoir quelle phase est au-dessus, il faut connaître les masses volumiques des deux solvants. Le liquide le moins dense se place au-dessus.
\end{correctionbox}

\subsection*{Correction exercice 9}
\begin{correctionbox}
\ce{CH3-CH2-OH} : famille des alcools, éthanol.

\ce{CH3-CHO} : famille des aldéhydes, éthanal.

\ce{CH3-CO-CH3} : famille des cétones, propanone.

\ce{CH3-COOH} : famille des acides carboxyliques, acide éthanoïque.
\end{correctionbox}

\subsection*{Correction exercice 10}
\begin{correctionbox}
Le rendement vaut :
\[
\eta=\frac{m_{\text{obtenu}}}{m_{\text{théorique}}}\times100
=\frac{6.20}{8.50}\times100=\SI{72.9}{\percent}.
\]
Le rendement est d'environ \SI{73}{\percent}. Il peut être inférieur à \SI{100}{\percent} à cause de pertes lors des manipulations, d'une transformation non totale, d'un isolement incomplet, ou d'une purification qui élimine une partie du produit.
\end{correctionbox}

\subsection*{Correction exercice 11}
\begin{correctionbox}
Méthane :
\[
\ce{CH4 + 2O2 -> CO2 + 2H2O}.
\]
Éthanol :
\[
\ce{C2H6O + 3O2 -> 2CO2 + 3H2O}.
\]
Butane :
\[
\ce{2C4H10 + 13O2 -> 8CO2 + 10H2O}.
\]
\end{correctionbox}

\subsection*{Correction exercice 12}
\begin{correctionbox}
Liaisons rompues : dans \ce{CH4}, il y a 4 liaisons \ce{C-H}. Dans $2\ce{O2}$, il y a 2 liaisons \ce{O=O}. Donc :
\[
E_{\text{rompues}}=4\times413+2\times498=1652+996=\SI{2648}{kJ.mol^{-1}}.
\]
Liaisons formées : dans \ce{CO2}, il y a 2 liaisons \ce{C=O}. Dans $2\ce{H2O}$, il y a 4 liaisons \ce{O-H}. Donc :
\[
E_{\text{formées}}=2\times799+4\times463=1598+1852=\SI{3450}{kJ.mol^{-1}}.
\]
Ainsi :
\[
\Delta_rE\approx2648-3450=\SI{-802}{kJ.mol^{-1}}.
\]
La combustion du méthane est exothermique.
\end{correctionbox}

\newpage

%====================================================
\section{Correction du devoir bilan final}

\subsection*{Correction de la partie A}
\begin{correctionbox}
Dans la réaction :
\[
\ce{MnO4^- + 5Fe^{2+} + 8H+ -> Mn^{2+} + 5Fe^{3+} + 4H2O},
\]
\ce{MnO4^-} est réduit en \ce{Mn^{2+}} : c'est l'oxydant. \ce{Fe^{2+}} est oxydé en \ce{Fe^{3+}} : c'est le réducteur.

À l'équivalence :
\[
\frac{n(\ce{Fe^{2+}})}{5}=n(\ce{MnO4^-}).
\]
Or :
\[
n(\ce{MnO4^-})=c_BV_E=1.00\times10^{-2}\times16.4\times10^{-3}=\SI{1.64e-4}{mol}.
\]
Donc :
\[
n(\ce{Fe^{2+}})=5\times1.64\times10^{-4}=\SI{8.20e-4}{mol}.
\]
Le volume titré vaut :
\[
V_A=\SI{20.0}{mL}=\SI{2.00e-2}{L}.
\]
La concentration est :
\[
c_A=\frac{8.20\times10^{-4}}{2.00\times10^{-2}}=\SI{4.10e-2}{mol.L^{-1}}.
\]
Avant l'équivalence, le permanganate versé est limitant ; après l'équivalence, les ions fer(II) sont consommés et le permanganate devient en excès. L'équivalence marque donc le changement de réactif limitant.
\end{correctionbox}

\subsection*{Correction de la partie B}
\begin{correctionbox}
L'éthanol possède le groupe hydroxyle \ce{-OH} : il appartient à la famille des alcools.

L'éthanol est miscible avec l'eau car son groupe \ce{-OH} peut établir des ponts hydrogène avec les molécules d'eau. Sa chaîne carbonée est courte, donc la partie hydrophile domine suffisamment.

\ce{CO2} possède des liaisons polarisées mais une géométrie linéaire symétrique : il est globalement apolaire. \ce{H2O} est coudée et ses liaisons \ce{O-H} sont polarisées : elle est polaire.

L'eau dissout de nombreux solides ioniques car ses molécules polaires stabilisent les ions séparés : les cations sont entourés par le côté oxygène $\delta^-$, les anions par les côtés hydrogène $\delta^+$.
\end{correctionbox}

\subsection*{Correction de la partie C}
\begin{correctionbox}
La masse maximale théorique vaut :
\[
m_{\max}=n_{\max}M=0.120\times88.0=\SI{10.56}{g}.
\]
Le rendement vaut :
\[
\eta=\frac{7.80}{10.56}\times100=\SI{73.9}{\percent}.
\]
Lors de la combustion complète d'un ester ne contenant que C, H et O, les produits formés sont \ce{CO2} et \ce{H2O}.

Risques : incendie, explosion, intoxication au monoxyde de carbone en cas de combustion incomplète. Enjeu environnemental : émissions de \ce{CO2}, contribution au changement climatique, nécessité de réduire l'empreinte carbone.
\end{correctionbox}

\newpage

%====================================================
\section{Fiche de synthèse finale}

\begin{retenirbox}
\textbf{Quantités de matière}
\[
n=\frac{m}{M},\qquad n=\frac{V}{V_m},\qquad c=\frac{n}{V},\qquad C_m=Mc.
\]
Toujours vérifier les unités : masse en grammes si $M$ est en \si{g.mol^{-1}}, volume en litres si $c$ est en \si{mol.L^{-1}}.
\end{retenirbox}

\begin{retenirbox}
\textbf{Beer-Lambert et étalonnage}
\[
A=kc.
\]
Une droite d'étalonnage permet de déterminer une concentration inconnue à partir d'une absorbance mesurée, dans le domaine de validité de la loi.
\end{retenirbox}

\begin{retenirbox}
\textbf{Oxydo-réduction}
Un oxydant capte des électrons ; un réducteur cède des électrons. Une réaction d'oxydo-réduction est un transfert d'électrons. Les demi-équations permettent d'établir l'équation globale.
\end{retenirbox}

\begin{retenirbox}
\textbf{Avancement}
Pour $aA+bB\rightarrow cC+dD$ :
\[
n_f(A)=n_i(A)-ax,\quad n_f(B)=n_i(B)-bx,
\]
\[
n_f(C)=n_i(C)+cx,\quad n_f(D)=n_i(D)+dx.
\]
L'avancement maximal est imposé par le réactif limitant.
\end{retenirbox}

\begin{retenirbox}
\textbf{Titrage}
À l'équivalence, les réactifs sont introduits dans les proportions stœchiométriques. Pour $aA+bB\rightarrow\cdots$ :
\[
\frac{n(A)}{a}=\frac{n(B)}{b}.
\]
\end{retenirbox}

\begin{retenirbox}
\textbf{Structure et propriétés}
Le schéma de Lewis donne accès à la géométrie. La géométrie et l'électronégativité permettent de déterminer la polarité. La polarité et les interactions expliquent la cohésion, la solubilité, la miscibilité et l'extraction.
\end{retenirbox}

\begin{retenirbox}
\textbf{Chimie organique}
Savoir reconnaître les familles : alcool, aldéhyde, cétone, acide carboxylique. Savoir exploiter un spectre IR, calculer un rendement et ajuster une combustion complète.
\end{retenirbox}

\begin{erreurbox}
\textbf{Erreurs à éviter absolument}
\begin{itemize}
\item oublier les coefficients stœchiométriques ;
\item confondre concentration en masse et concentration molaire ;
\item utiliser des mL au lieu de L dans $n=cV$ ;
\item oublier qu'à l'équivalence les réactifs sont en proportions stœchiométriques ;
\item conclure sur la polarité d'une molécule sans tenir compte de sa géométrie ;
\item confondre dissolution et réaction chimique de transformation ;
\item calculer un rendement avec une masse théorique correspondant au mauvais réactif ;
\item oublier que rompre des liaisons consomme de l'énergie et former des liaisons en libère.
\end{itemize}
\end{erreurbox}

\begin{center}
\Large\bfseries Fin du cours
\end{center}

\end{document}
